%% ==============================================================================
%% frontmatter/acknowledgements.tex
%% Remerciements et hommage aux artisans de l'écosystème TeX
%% ==============================================================================
\addchap{Remerciements et Avant-Propos}

Ce manuel a été conçu avec la volonté de démystifier l'architecture des grands projets \LaTeX{}
et d'offrir un modèle de départ pérenne, moderne et immédiatement exploitable.

Mes remerciements les plus sincères vont d'abord aux pionniers de l'informatique typographique :
à Donald Knuth pour avoir imaginé un moteur de composition d'une précision inégalée,
à Leslie Lamport pour avoir rendu ce système accessible aux scientifiques,
ainsi qu'à Markus Kohm, dont la suite KOMA-Script transcende les limites des classes historiques
en apportant la rigueur et l'élégance de la tradition éditoriale européenne.

J'exprime toute ma reconnaissance à l'équipe du projet \LaTeX3,
en particulier à Frank Mittelbach et Ulrike Fischer,
pour leur travail monumental de modernisation du noyau \LaTeX{} contemporain
et l'intégration harmonieuse des nouvelles couches de programmation et de gestion des sockets.

Merci également aux concepteurs de Bib\LaTeX{} et de Biber,
qui ont su libérer la bibliographie académique des entraves de l'ancien format BST,
ouvrant la voie à une personnalisation puissante et transparente.

Puisse ce tutoriel vous faire gagner un temps précieux dans la structuration de vos mémoires,
thèses, rapports d'ingénierie et monographies scientifiques.
